\section{Informatik}

\subsection{Relationale vs. multidimensionale Datenbanken}

\definition{Eine \textbf{Datenbank} (DB) ist eine strukturierte, persistente Sammlung von Daten. Ein \textbf{Datenbankmanagementsystem} (DBMS) ist die Software zur Verwaltung von Datenbanken (z.\,B.\  PostgreSQL, Oracle, MySQL).}

\textbf{OLTP vs. OLAP -- Verarbeitungsparadigmen:}
\begin{itemize}[leftmargin=1.5em]
    \item \textbf{OLTP} (Online Transaction Processing): Operative Verarbeitung mit kleinen Datenmengen pro Transaktion, hohe Schreiblast (INSERT, UPDATE, DELETE), Antwortzeit Millisekunden. Beispiele: Bank, Webshop.
    \item \textbf{OLAP} (Online Analytical Processing): Analytische Auswertungen mit sehr großen Datenvolumina, komplexe SELECT-Queries mit Aggregationen, Antwortzeit Sekunden bis Minuten. Beispiele: Business Intelligence, Reporting, Data Mining.
\end{itemize}

\textbf{Datenmodelle:}\vspace{2pt}
\begin{itemize}[leftmargin=1.5em]
    \item \textbf{OLTP:} Relational, normalisiert (3NF, BCNF) -- Redundanzreduktion, Schreiboptimierung
    \item \textbf{OLAP:} Multidimensional, denormalisiert -- bewusste Redundanz für Lesezugriffe
\end{itemize}

\definition{\textbf{Data Warehouse} (DWH) ist eine zentrale, integrierte, historisierte Datenbasis, optimiert für analytische Abfragen. Daten werden aus OLTP-Systemen via \textbf{ETL-Prozess} zugeführt (Extract, Transform, Load).}

\textbf{Schema-Formen im DWH:}\vspace{2pt}
\begin{itemize}[leftmargin=1.5em]
    \item \textbf{Star-Schema:} Zentrale \textbf{Faktentabelle} mit Kennzahlen (Measures) + direkt angebundene \textbf{Dimensionstabellen} (Zeit, Produkt, Region, Kunde)
    \item \textbf{Snowflake-Schema:} Normalisierte Dimensionstabellen mit Unterdimensionen; komplexer, aber speichereffizienter
\end{itemize}

\textbf{OLAP-Würfel und Hierarchien:}\vspace{2pt}
\begin{itemize}[leftmargin=1.5em]
    \item \textbf{Measures (Fakten):} Umsatz, Gewinn, Menge
    \item \textbf{Dimensionen:} Zeit, Produkt, Kunde, Region
    \item \textbf{Hierarchien:} Jahr $\to$ Quartal $\to$ Monat; Kontinent $\to$ Land $\to$ Stadt
\end{itemize}

\textbf{Typische OLAP-Operationen:}\vspace{2pt}
\begin{itemize}[leftmargin=1.5em]
    \item \textbf{Slice:} Eine Dimension fixieren (z.\,B.\  nur 2023)
    \item \textbf{Dice:} Teilwürfel auswahl (z.\,B.\  2023 und Region='Europa')
    \item \textbf{Drill-Down:} Detailliertere hierarchische Ebene
    \item \textbf{Drill-Up/Roll-Up:} Aggregation auf höhere Ebene
\end{itemize}

\textbf{OLAP-Varianten:}\vspace{2pt}
\begin{itemize}[leftmargin=1.5em]
    \item \textbf{MOLAP:} Multidimensionale Speicherung in Speicherkuben; schnell, aber Speicher
    \item \textbf{ROLAP:} Relationale Speicherung; flexibel, aber langsamer
    \item \textbf{HOLAP:} Hybridansatz; Kombination beider
\end{itemize}

\definition{\textbf{Data Lake} ist eine Speicherung von Rohdaten in beliebigem Format (strukturiert, semi-strukturiert, unstrukturiert). Flexibler als DWH, aber höhere Anforderungen an Governance und Datenqualität.}

\newpage

\subsection{Entity-Relationship-Modell}

\definition{Das \textbf{Entity-Relationship-Modell} (ERM) ist ein konzeptionelles Datenmodellierungsverfahren zur fachlichen Abbildung der Miniwelt. Es besteht aus \textbf{Entities} (Objekten), \textbf{Attributen} (Eigenschaften) und \textbf{Relationships} (Beziehungen).}

\textbf{Datenmodellierungsebenen:}\vspace{2pt}
\begin{itemize}[leftmargin=1.5em]
    \item \textbf{Konzeptionell:} Fachliche Sicht (ER-Modell)
    \item \textbf{Logisch:} Relationales Schema (Tabellen, Spalten)
    \item \textbf{Physisch:} Konkrete Speicherung im DBMS
\end{itemize}

\textbf{Schlüsselbegriffe:}\vspace{2pt}
\begin{itemize}[leftmargin=1.5em]
    \item \textbf{Superkey:} Attributkombination, die Tupel eindeutig identifiziert
    \item \textbf{Candidate Key:} Minimaler Superkey
    \item \textbf{Primärschlüssel (PK):} Gewählter Candidate Key
    \item \textbf{Fremdschlüssel (FK):} Verweist auf PK einer anderen Tabelle
    \item \textbf{Surrogate Key:} Künstlicher Schlüssel (SERIAL, AUTO\_INCREMENT)
\end{itemize}

\textbf{Kardinalitäten und Beziehungstypen:}\vspace{2pt}
\begin{itemize}[leftmargin=1.5em]
    \item \textbf{1:1} (Eins-zu-eins): Person $\leftrightarrow$ Reisepass
    \item \textbf{1:N} (Eins-zu-viele): Abteilung $\leftrightarrow$ Mitarbeiter
    \item \textbf{N:M} (Viele-zu-viele): Student $\leftrightarrow$ Kurs
\end{itemize}

\textbf{Optionalität / Teilnahmebedingungen:}\vspace{2pt}
\begin{itemize}[leftmargin=1.5em]
    \item \textbf{Optional:} 0..1 oder 0..N (Teilnahme nicht erforderlich)
    \item \textbf{Verpflichtend:} 1..1 oder 1..N (Teilnahme erforderlich)
\end{itemize}

\definition{Eine \textbf{Schwache Entität} besitzt keinen vollständigen eigenen Schlüssel und wird über eine identifizierende Beziehung zu einer starken Entität bestimmt. Beispiel: Bestellposition ist schwach abhängig von Bestellung.}

\textbf{Relationales Modell -- Grundkonzepte:}\vspace{2pt}
\begin{itemize}[leftmargin=1.5em]
    \item \textbf{Relation:} Tabelle
    \item \textbf{Tupel:} Zeile (Record)
    \item \textbf{Attribut:} Spalte
    \item \textbf{Domäne:} Wertebereich eines Attributs
\end{itemize}

\textbf{Integritätsbedingungen:}\vspace{2pt}
\begin{itemize}[leftmargin=1.5em]
    \item \textbf{Entitätsintegrität:} PK eindeutig und NOT NULL
    \item \textbf{Referentielle Integrität:} FK muss existierenden PK referenzieren
    \item \textbf{Domänenintegrität:} Werte innerhalb erlaubter Domäne
\end{itemize}

\textbf{Abbildung ERM $\to$ Relationales Modell:}\vspace{2pt}
\begin{itemize}[leftmargin=1.5em]
    \item \textbf{Entity $\to$ Tabelle:} Jede Entity-Menge wird zu einer Tabelle
    \item \textbf{Attribute $\to$ Spalten:} Attribute werden zu Spalten
    \item \textbf{1:N $\to$ FK:} FK auf N-Seite einführen
    \item \textbf{N:M $\to$ Verbindungstabelle:} Neue Tabelle mit FKs zu beiden Entitäten
    \item \textbf{1:1 $\to$ FK oder Zusammenführung:} Entweder FK oder Relationstabellen zusammenführen
\end{itemize}

\newpage

\subsection{SQL als Data Query Language}

\definition{\textbf{SQL} (Structured Query Language) ist die standardisierte Abfragesprache für relationale Datenbanksysteme. Sie wird in Kategorien eingeteilt: DQL (Abfragen), DML (Datenänderungen), DDL (Definitionen), DCL (Kontrolle), TCL (Transaktionen).}

\textbf{SELECT-Abfrage -- Grundstruktur:}\vspace{2pt}
\begin{itemize}[leftmargin=1.5em]
    \item \textbf{FROM:} Quellentabelle(n)
    \item \textbf{WHERE:} Zeilenfilterung (Selektion)
    \item \textbf{GROUP BY:} Gruppierung nach Attributen
    \item \textbf{HAVING:} Filterung auf Gruppen
    \item \textbf{SELECT:} Spaltenprojektion und Transformationen
    \item \textbf{ORDER BY:} Sortierung
    \item \textbf{LIMIT/OFFSET:} Paginierung
\end{itemize}

\textbf{Logische Ausführungsreihenfolge:}\vspace{2pt}
\begin{enumerate}[leftmargin=1.5em]
    \item FROM $\to$ Tabelle laden
    \item WHERE $\to$ Zeilenfilterung
    \item GROUP BY $\to$ Gruppierung
    \item HAVING $\to$ Gruppenfilterung
    \item SELECT $\to$ Spaltenprojektion
    \item ORDER BY $\to$ Sortierung
    \item LIMIT/OFFSET $\to$ Paginierung
\end{enumerate}

\definition{\textbf{NULL-Semantik:} NULL bedeutet "unbekannt", nicht Null oder leerer String. Vergleiche mit NULL-Werten geben UNKNOWN zurück. Stattdessen: \texttt{IS NULL} oder \texttt{IS NOT NULL}.}

\textbf{Wichtige WHERE-Operatoren:}\vspace{2pt}
\begin{itemize}[leftmargin=1.5em]
    \item \textbf{Vergleich:} =, <>, >, <, >=, <=
    \item \textbf{BETWEEN:} z.\,B.\  \texttt{preis BETWEEN 10 AND 100}
    \item \textbf{IN:} z.\,B.\  \texttt{stadt IN ('Wien', 'Graz')}
    \item \textbf{LIKE:} Mustervergleich mit \texttt{\_} (1 Zeichen) und \texttt{\%} (beliebig)
    \item \textbf{IS NULL / IS NOT NULL:} Null-Vergleiche
\end{itemize}

\textbf{Aggregatfunktionen:}\vspace{2pt}
\begin{itemize}[leftmargin=1.5em]
    \item \textbf{COUNT(*):} Zahl aller Zeilen
    \item \textbf{SUM(spalte):} Summe
    \item \textbf{AVG(spalte):} Durchschnitt
    \item \textbf{MIN(spalte), MAX(spalte):} Minimum, Maximum
\end{itemize}

\textbf{JOINs -- Tabellenkombination:}\vspace{2pt}
\begin{itemize}[leftmargin=1.5em]
    \item \textbf{INNER JOIN:} Nur Zeilen mit Match in beiden Tabellen (Schnittmenge)
    \item \textbf{LEFT JOIN:} Alle Zeilen der linken Tabelle + Matches rechts (NULL für fehlende)
    \item \textbf{RIGHT JOIN:} Alle Zeilen der rechten Tabelle + Matches links
    \item \textbf{FULL OUTER JOIN:} Alle Zeilen aus beiden Tabellen
    \item \textbf{CROSS JOIN:} Kartesisches Produkt (alle Kombinationen)
\end{itemize}

\textbf{Weitere wichtige Konzepte:}\vspace{2pt}
\begin{itemize}[leftmargin=1.5em]
    \item \textbf{DISTINCT:} Entfernt Duplikate in Ergebnissen
    \item \textbf{Subqueries:} Verschachtelte SELECT-Abfragen in WHERE, FROM oder SELECT
    \item \textbf{Views:} Virtuelle Tabellen auf Basis von Abfragen
    \item \textbf{Indizes:} Beschleunigen Suchzugriffe auf Spalten
    \item \textbf{COALESCE(a,b,...):} Erster Nicht-NULL-Wert
\end{itemize}

\newpage

\subsection{Normalformen, Integritätsbedingungen und Transaktionen}

\definition{\textbf{Normalisierung} ist das systematische Umstrukturieren eines Relationenschemata mit dem Ziel, Redundanzen zu reduzieren und Anomalien zu vermeiden.}

\textbf{Anomalien (Probleme ohne Normalisierung):}\vspace{2pt}
\begin{itemize}[leftmargin=1.5em]
    \item \textbf{Insert-Anomalie:} Neue Entität kann nicht eingefügt werden ohne zusätzliche irrelevante Daten
    \item \textbf{Update-Anomalie:} Informationen müssen an mehreren Stellen aktualisiert werden (Inkonsistenzrisiko)
    \item \textbf{Delete-Anomalie:} Löschen führt zu ungewolltem Datenverlust
\end{itemize}

\definition{\textbf{1. Normalform (1NF):} Alle Attributwerte sind atomar (unteilbar); keine Wiederholungsgruppen.}

\definition{\textbf{2. Normalform (2NF):} Erfüllt 1NF und es gibt keine partiellen Abhängigkeiten (relevant bei zusammengesetztem PK).}

\definition{\textbf{3. Normalform (3NF):} Erfüllt 2NF und es gibt keine transitiven Abhängigkeiten (kein Nicht-Schlüsselattribut hängt von anderem ab).}

\definition{\textbf{Boyce-Codd-Normalform (BCNF):} Für jede funktionale Abhängigkeit $X \to Y$ muss $X$ ein Superschlüssel sein. Strenger als 3NF, eliminiert weitere Spezialfälle.}

\textbf{Integritätsbedingungen:}\vspace{2pt}
\begin{itemize}[leftmargin=1.5em]
    \item \textbf{Entitätsintegrität:} PK ist eindeutig und nicht NULL
    \item \textbf{Referentielle Integrität:} FK muss auf existierenden PK referenzieren
    \item \textbf{Domänenintegrität:} Werte müssen in erlaubter Domäne liegen
\end{itemize}

\textbf{Constraints (SQL-Implementierung):}\vspace{2pt}
\begin{itemize}[leftmargin=1.5em]
    \item \textbf{PRIMARY KEY:} Eindeutigkeits- und NOT-NULL-Constraint
    \item \textbf{FOREIGN KEY:} Referentielle Integrität
    \item \textbf{UNIQUE:} Eindeutigkeit (NULL mehrmals erlaubt)
    \item \textbf{NOT NULL:} Wert erforderlich
    \item \textbf{CHECK:} Bedingung auf Wert
    \item \textbf{DEFAULT:} Standardwert
\end{itemize}

\textbf{Referenzielle Aktionen (ON DELETE / ON UPDATE):}\vspace{2pt}
\begin{itemize}[leftmargin=1.5em]
    \item \textbf{CASCADE:} Automatisch gelöscht/aktualisieren
    \item \textbf{SET NULL:} FK auf NULL setzen
    \item \textbf{RESTRICT:} Verweigern falls FK existiert
\end{itemize}

\definition{Eine \textbf{Transaktion} ist eine logische Einheit mehrerer Datenbankoperationen. Alle Operationen werden vollständig (COMMIT) oder gar nicht (ROLLBACK) ausgeführt.}

\definition{\textbf{ACID-Eigenschaften:}
\begin{itemize}[leftmargin=1.5em]
    \item \textbf{Atomicity:} Alles oder Nichts
    \item \textbf{Consistency:} Datenbankbleibt in konsistentem Zustand
    \item \textbf{Isolation:} Parallele Transaktionen beeinflussen sich nicht
    \item \textbf{Durability:} Ge committete Daten persistent
\end{itemize}

\textbf{Isolationslevel (gegen Anomalien):}\vspace{2pt}
\begin{itemize}[leftmargin=1.5em]
    \item \textbf{Read Uncommitted:} Niedrigster Standard; Dirty Reads möglich
    \item \textbf{Read Committed:} Verhindert Dirty Reads; aber Non-repeatable Reads möglich
    \item \textbf{Repeatable Read:} Verhindert Non-repeatable Reads; aber Phantom Reads möglich
    \item \textbf{Serializable:} Höchste Isolation; verhindert alle Anomalien
\end{itemize}

\definition{\textbf{Write-Ahead Logging (WAL):} Änderungen werden zuerst ins Log geschrieben und danach in die Datenbank. Wichtig für Durability und Recovery nach Ausfällen.}

\newpage

\subsection{Datenschutz und Datensicherheit}

\subsubsection{Begriffsabgrenzung}

\textbf{Datenschutz} und \textbf{Datensicherheit} werden im Alltag oft synonym verwendet, bezeichnen aber unterschiedliche Konzepte:

\begin{itemize}
  \item \textbf{Datenschutz} (engl. \textit{data privacy}) ist das \textbf{rechtliche und organisatorische} Konzept, das den Schutz personenbezogener Daten natürlicher Personen vor unbefugter Verarbeitung sicherstellt. Es geht primär um die Frage: \emph{Darf} ein Unternehmen diese Daten überhaupt erheben und verarbeiten?
  \item \textbf{Datensicherheit} (engl. \textit{data security / information security}) ist das \textbf{technische und organisatorische} Konzept, das alle Daten -- unabhängig ob personenbezogen -- vor Verlust, Manipulation, unbefugtem Zugriff und anderen Bedrohungen schützt. Es geht um die Frage: \emph{Wie} werden Daten technisch gesichert?
\end{itemize}

\subsubsection{Rechtlicher Rahmen: DSGVO}

Die \textbf{Datenschutz-Grundverordnung} (DSGVO), in Kraft seit Mai 2018, ist die maßgebliche europäische Rechtsgrundlage für den Umgang mit personenbezogenen Daten. Sie gilt für alle Unternehmen, die Daten von EU-Bürgern verarbeiten -- unabhängig vom Unternehmensstandort.

\textbf{Grundsätze der DSGVO:}
\begin{enumerate}
  \item \textbf{Rechtmäßigkeit, Transparenz:} Verarbeitung braucht eine Rechtsgrundlage (Einwilligung, Vertrag, rechtliche Verpflichtung, berechtigte Interessen)
  \item \textbf{Zweckbindung:} Daten dürfen nur für den bei der Erhebung festgelegten Zweck verwendet werden
  \item \textbf{Datenminimierung:} Es dürfen nur die Daten erhoben werden, die für den Zweck erforderlich sind
  \item \textbf{Richtigkeit:} Daten müssen korrekt und aktuell sein
  \item \textbf{Speicherbegrenzung:} Speicherung nur solange nötig
  \item \textbf{Integrität und Vertraulichkeit:} Datenschutz durch Maßnahmen
  \item \textbf{Rechenschaftspflicht:} Der Verantwortliche muss Einhaltung nachweisen
\end{enumerate}

\textbf{Betroffenenrechte:} Recht auf Auskunft, Berichtigung, Löschung (``Recht auf Vergessenwerden''), Datenportabilität, Widerspruch gegen automatisierte Entscheidungen.

\textbf{Sanktionen:} Verstöße können mit bis zu \textbf{20 Millionen Euro} oder \textbf{4\,\% des weltweiten Jahresumsatzes} geahndet werden.

\subsubsection{CIA-Modell der Informationssicherheit}

Das klassische Rahmenwerk definiert drei Schutzziele:

\begin{itemize}
  \item \textbf{Confidentiality (Vertraulichkeit):} Unbefugte dürfen nicht auf Daten zugreifen. Maßnahmen: Verschlüsselung, Zugriffskontrolle
  \item \textbf{Integrity (Integrität):} Daten dürfen nicht unbefugt verändert werden. Maßnahmen: Hashwerte, digitale Signaturen
  \item \textbf{Availability (Verfügbarkeit):} Autorisierte Benutzer sollen jederzeit auf Daten zugreifen können. Maßnahmen: Redundanz, Backups, DDoS-Schutz
\end{itemize}

\subsubsection{Technische und organisatorische Maßnahmen}

Die DSGVO verlangt ``Privacy by Design'' sowie angemessene TOMs:

\textbf{Wichtige Maßnahmen:}
\begin{itemize}
  \item Pseudonymisierung und Anonymisierung
  \item Verschlüsselung (AES-256 für Daten, TLS 1.3 für Übertragung)
  \item Role-Based Access Control (RBAC)
  \item Backups nach 3-2-1-Regel (3 Kopien, 2 Medientypen, 1 Off-site)
  \item Umfassendes Logging und Auditing
  \item Regelmäßige Schulungen der Mitarbeiter
\end{itemize}

\subsubsection{Datenschutz im ML/Data-Science-Kontext}

In Data-Science-Projekten ergeben sich spezifische Herausforderungen:

\begin{itemize}
  \item \textbf{Datenvergiftung:} Unvorsichtige Nutzung von Trainingsdaten kann zu Datenlecks führen
  \item \textbf{Model Inversion:} ML-Modelle können trainierte Daten teilweise rekonstruieren
  \item \textbf{Fairness und Bias:} Diskriminierende Modelle verletzen DSGVO
  \item \textbf{Datentransfers:} Übermittlung von EU-Daten in Drittländer ist oft nicht zulässig
\end{itemize}

\newpage

\subsection{Python: Listen, Tupel und Mengen}

\definition{\textbf{Listen} (\texttt{list}), \textbf{Tupel} (\texttt{tuple}) und \textbf{Mengen} (\texttt{set}) sind fundamentale Sequenz- bzw.\ Sammlungstypen in Python mit unterschiedlichen semantischen Garantien und Laufzeiteigenschaften.}

\textbf{Gemeinsamkeiten:}\vspace{2pt}
\begin{itemize}[leftmargin=1.5em]
  \item Alle sind \textbf{iterierbar} -- durchlaufbar in \texttt{for}-Schleife
  \item Unterstützen \texttt{in}-Operator und Funktionen \texttt{len()}, \texttt{min()}, \texttt{max()}, \texttt{sum()}
  \item Können heterogene Python-Objekte aufnehmen
\end{itemize}

\textbf{Die Liste -- veränderliche Sequenz:}\vspace{2pt}
\begin{itemize}[leftmargin=1.5em]
  \item \textbf{Geordnet} (Einfügereihenfolge bleibt erhalten), \textbf{veränderlich} (mutable)
  \item Index-Zugriff $\mathcal{O}(1)$, Einfügen an beliebiger Position $\mathcal{O}(n)$
  \item Einsatz: Dynamische Sequenzen mit Laufzeitänderungen
\end{itemize}

\textbf{Das Tupel -- unveränderliche Sequenz:}\vspace{2pt}
\begin{itemize}[leftmargin=1.5em]
  \item \textbf{Geordnet}, aber \textbf{unveränderlich} (immutable)
  \item \textbf{Hashbar}: Kann als Dictionary-Key oder in \texttt{set} verwendet werden
  \item Speichereffizienter, kein Größen-Overhead wie bei Listen
  \item Einsatz: Datenbankdatensätze, Funktionsrückgaben, Konfigurationsparameter
  \item Named Tuples: Tupel mit benannten Feldern via \texttt{typing.NamedTuple}
\end{itemize}

\textbf{Die Menge -- ungeordnete, eindeutige Kollektion:}\vspace{2pt}
\begin{itemize}[leftmargin=1.5em]
  \item \textbf{Ungeordnet}, Duplikate automatisch verworfen
  \item Intern Hash-Tabelle: Mitgliedschaftstest, Einfügen, Löschen $\mathcal{O}(1)$ Ø
  \item Unterstützt Mengenoperationen: Vereinigung (\texttt{|}), Schnitt (\texttt{\&}), Differenz (\texttt{-}), sym. Differenz (\texttt{\^})
  \item \textbf{frozenset}: Unveränderliches Pendant (hashbar)
  \item Einsatz: Schnelle Eindeutigkeitstests, Deduplizierung, Mengenlogik
\end{itemize}

\begin{table}[h!]
\centering
\rowcolors{2}{tablerowB}{tablerowA}
\begin{tabularx}{\linewidth}{L{2.8cm} C{2cm} C{2cm} C{3.2cm}}
\rowcolor{tablehead}
\color{white}\textbf{Eigenschaft} &
\color{white}\textbf{list} &
\color{white}\textbf{tuple} &
\color{white}\textbf{set} \\
\midrule
Geordnet & \checkmark & \checkmark & -- \\
Veränderlich & mutable & immutable & mutable \\
Duplikate & \checkmark & \checkmark & -- \\
Mitgliedschaftstest & $\mathcal{O}(n)$ & $\mathcal{O}(n)$ & $\mathcal{O}(1)$ \\
Hashbar & -- & \checkmark & -- \\
\end{tabularx}
\end{table}

\newpage

\subsection{Python: Kopien vs. Verweise, Shallow \& Deep Copies}

\definition{In Python ist jede Variable ein \textbf{Name (Label)} auf ein Objekt im Speicher. Eine Zuweisung \texttt{b = a} kopiert nicht das Objekt, sondern lässt beide auf \textbf{dasselbe} Objekt zeigen.}

\textbf{Operator \texttt{is} vs.\ \texttt{==}:}\vspace{2pt}
\begin{itemize}[leftmargin=1.5em]
  \item \textbf{\texttt{is}}: Prüft Identität (gleiche Speicheradresse, Pointer-Vergleich)
  \item \textbf{\texttt{==}}: Prüft Gleichheit (gleicher Wert)
  \item Bei unveränderlichen Typen praktisch egal; bei mutable Typen kritisch
\end{itemize}

\textbf{Immutable vs. Mutable:}\vspace{2pt}
\begin{itemize}[leftmargin=1.5em]
  \item \textbf{Immutable}: \texttt{int, float, str, tuple, frozenset} -- scheinbare Änderungen erzeugen neue Objekte
  \item \textbf{Mutable}: \texttt{list, dict, set} -- Attribute/Elemente können in-place verändert werden
\end{itemize}

\definition{\textbf{Shallow Copy}}: Neuer Container auf oberster Ebene, aber enthaltene Elemente werden nicht kopiert. Bei verschachtelten Strukturen teilen Original und Kopie die inneren Objekte.}

\textbf{Wege zur Shallow Copy:}\vspace{2pt}
\begin{itemize}[leftmargin=1.5em]
  \item \texttt{list(orig)}, \texttt{orig[:]}, \texttt{copy.copy(orig)}, \texttt{orig.copy()}
\end{itemize}

\definition{\textbf{Deep Copy}: Vollständig unabhängiges Objekt inklusive aller rekursiv erreichbaren Unterobjekte. \textbf{Keine} gemeinsamen mutable Objekte. Behandelt Zyklen via internes Memo-Dict.}

\textbf{Performance-Aspekt:}\vspace{2pt}
\begin{itemize}[leftmargin=1.5em]
  \item Deep Copy ist deutlich teurer -- gesamter Objektgraph wird traversiert
  \item Für Performance-kritischen Code vermeiden; stattdessen unveränderliche Strukturen oder explizite Rekonstruktion
\end{itemize}

\textbf{Funktionen und Parameterübergabe (Pass by Object Reference):}\vspace{2pt}
\begin{itemize}[leftmargin=1.5em]
  \item Funktion erhält Referenz auf Objekt; kann mutable Objekte in-place verändern
  \item Neuzuweisung des Parameters beeinflusst aufrufende Seite nicht
\end{itemize}

\begin{merksatz}[Default-Argument-Falle]
\texttt{def f(lst=[]): ...} -- leere Liste wird einmal erstellt und über alle Aufrufe geteilt!\\Korrekt: \texttt{def f(lst=None): if lst is None: lst = []}
\end{merksatz}

\newpage

\subsection{Python: Schritte des Data-Preprocessings (scikit-learn)}

\definition{\textbf{Data-Preprocessing} ist die systematische Vorbereitung von Rohdaten für ML-Algorithmen. Viele Verfahren (SVM, k-NN, neuronale Netze) sind empfindlich gegenüber Skalierungen, fehlenden Werten oder falschen Kodierungen.}

\textbf{Scikit-learn Pipeline-Konzept:}\vspace{2pt}
\begin{itemize}[leftmargin=1.5em]
  \item Jeder Schritt hat \texttt{fit(X)} (auf Trainingsdaten lernen) und \texttt{transform(X)} (anwenden)
  \item Verkettung via \texttt{Pipeline}; kritisch: \texttt{fit} nur auf Trainingsdaten! (verhindert Data Leakage)
\end{itemize}

\textbf{Schritt 1 -- Fehlende Werte (Imputation):}\vspace{2pt}
\begin{itemize}[leftmargin=1.5em]
  \item \texttt{SimpleImputer}: Strategien -- \texttt{'mean'}, \texttt{'median'}, \texttt{'most\_frequent'}, \texttt{'constant'}
  \item \texttt{'median'} robust gegenüber Ausreißern; für kategorische: \texttt{'most\_frequent'}
  \item \texttt{KNNImputer}: Multivariate Imputation mit anderen Features
\end{itemize}

\textbf{Schritt 2 -- Feature-Skalierung:}\vspace{2pt}
\begin{itemize}[leftmargin=1.5em]
  \item \textbf{StandardScaler}: Z-Score $z = (x-\mu)/\sigma$ -- $\mu=0, \sigma=1$
  \item \textbf{MinMaxScaler}: Normalisierung auf $[0,1]$ -- empfindlich gegen Ausreißer
  \item \textbf{RobustScaler}: IQR-basiert $(x-\mathrm{Median})/\mathrm{IQR}$ -- ausreißer-robust
\end{itemize}

\textbf{Schritt 3 -- Kodierung kategorischer Features:}\vspace{2pt}
\begin{itemize}[leftmargin=1.5em]
  \item \textbf{OrdinalEncoder}: Ganze Zahlen (nur bei echten Ordnungen!)
  \item \textbf{OneHotEncoder}: Binäre Dummy-Variablen (keine implizite Rangfolge, aber hochdimensional)
\end{itemize}

\textbf{Schritt 4 -- Feature Engineering:}\vspace{2pt}
\begin{itemize}[leftmargin=1.5em]
  \item \textbf{PolynomialFeatures}: Interaktionsterme und Potenzen ($x_1, x_2 \to x_1^2, x_1 x_2, x_2^2$)
  \item \textbf{PowerTransformer}: Macht schiefe Verteilungen annähernd normal (Yeo-Johnson)
\end{itemize}

\textbf{Schritt 5 -- Feature-Selektion:}\vspace{2pt}
\begin{itemize}[leftmargin=1.5em]
  \item \texttt{VarianceThreshold}: Entfernt Features mit geringer Varianz
  \item \texttt{SelectKBest}: Top-$k$ Features nach Teststatistik
  \item \texttt{SelectFromModel}: Feature-Importances eines ersten Modells
\end{itemize}

\textbf{Schritt 6 -- Class Imbalance:}\vspace{2pt}
\begin{itemize}[leftmargin=1.5em]
  \item Stark unbalancierte Klassen $\to$ Modell klassifiziert alles als Majoritätsklasse
  \item Abhilfe: SMOTE (Oversampling), oder \texttt{class\_weight='balanced'}
\end{itemize}

\definition{\textbf{ColumnTransformer}: Wendet verschiedene Pipelines parallel auf unterschiedliche Feature-Typen an (numerisch, kategorisch).}

\newpage

\subsection{NumPy -- Einsatzmöglichkeiten und Vorteile gegenüber \glqq plain\grqq{} Python}

\definition{\textbf{NumPy} (Numerical Python) ist die Fundamentalbibliothek für numerisches Rechnen. Sie stellt homogene, multidimensionale Arrays (\texttt{numpy.ndarray}) und mathematische Funktionen bereit. Basis für SciPy, pandas, scikit-learn, PyTorch, TensorFlow.}

\textbf{Warum NumPy schnell ist:}\vspace{2pt}
\begin{itemize}[leftmargin=1.5em]
  \item \textbf{Homogener, kompakter Speicher}: Alle Elemente als C-Typen (z.B. \texttt{float64} = 8 Byte) in zusammenhängendem Block
  \item \textbf{Vektorisierung}: Operationen ohne explizite Python-Schleife -- nutzen SIMD (SSE, AVX)
  \item \textbf{Broadcasting}: Arrays unterschiedlicher Formen kombinierbar ohne Datenkopie
  \item \textbf{C/Fortran}: Kernroutinen in C/Fortran, nicht Python
  \item \textbf{Performance}: Tyischerweise 50-100x schneller als plain Python
\end{itemize}

\textbf{Kernkonzepte des ndarray:}\vspace{2pt}
\begin{itemize}[leftmargin=1.5em]
  \item \textbf{dtype}: Einheitlicher Datentyp (\texttt{float64, int32, bool, complex128})
  \item \textbf{Shape}: Dimensionen, z.B. \texttt{(100, 28, 28)} für 100 Bilder
  \item \textbf{Views vs. Copies}: Viele Ops (Slicing, Transpose, Reshape) geben View zurück -- selber Speicher! Mit \texttt{arr.base} prüfbar
\end{itemize}

\textbf{Broadcasting-Regeln:}\vspace{2pt}
\begin{itemize}[leftmargin=1.5em]
  \item Dimensionen werden von rechts ausgerichtet
  \item Eine Dimension der Größe 1 wird konzeptuell auf Größe des anderen Arrays ausgedehnt
  \item Keine physische Datenkopie -- konzeptuell
\end{itemize}

\textbf{Wichtige NumPy-Operationen:}\vspace{2pt}
\begin{itemize}[leftmargin=1.5em]
  \item \textbf{Erstellung}: \texttt{np.zeros, np.ones, np.arange, np.linspace, np.random.rand}
  \item \textbf{Indexing}: \texttt{a[2:5]}, \texttt{a[::2]}, \texttt{a[a > 0]} (Boolean Mask)
  \item \textbf{Aggregation}: \texttt{sum(), mean(), std(), max(axis=0)}
  \item \textbf{Lineare Algebra}: \texttt{np.dot}, \texttt{a @ b}, \texttt{np.linalg.inv, eig, solve}
  \item \textbf{Ufuncs}: \texttt{np.exp, log, sqrt, sin} -- elementweise, vektorisiert
\end{itemize}

\textbf{Einsatzmöglichkeiten:}\vspace{2pt}
\begin{itemize}[leftmargin=1.5em]
  \item \textbf{Lineare Algebra}: Matrixmultiplikation, Gleichungssysteme, Eigenwertzerlegung, FFT
  \item \textbf{Machine Learning}: Feature-Matrizen (scikit-learn), kompatibel mit PyTorch-Tensoren
  \item \textbf{Bilddaten}: RGB-Bild = \texttt{(H, W, 3)}-Array, Normalisierung, Cropping
  \item \textbf{Simulation}: ODE-Schritte, Wetterdaten-Interpolation, ohne Python-Schleifen
\end{itemize}

\textbf{Grenzen von NumPy:}\vspace{2pt}
\begin{itemize}[leftmargin=1.5em]
  \item Nur CPU; für GPU: CuPy, PyTorch, TensorFlow
  \item Keine automatische Differentiation; PyTorch oder JAX
\end{itemize}

\newpage

\subsection{R: Vektoren, Matrizen, Arrays, DataFrames, Listen -- Verwendungsmöglichkeiten sowie Vor- und Nachteile}

\definition{\textbf{R} ist eine funktionale, objektorientierte Sprache für Statistik und Data Science. Die Grundphilosophie: Alles ist ein Vektor (\texttt{atomare} Vektoren, Listen sind generische Vektoren).}

\textbf{Vektoren -- Grundelement:}\vspace{2pt}
\begin{itemize}[leftmargin=1.5em]
  \item \textbf{Atomar}: \texttt{numeric, integer, logical, character, complex, raw} -- homogen, geordnet
  \item Vektorisierte Operationen: \texttt{x + y, sin(x)} ohne Schleife
  \item Recycling: Kurzer Vektor wird wiederholt, bis Länge passt
  \item Einsatz: Numerische Berechnungen, Zeitreihen, Messreihen
\end{itemize}

\textbf{Matrizen -- 2D-Vektoren:}\vspace{2pt}
\begin{itemize}[leftmargin=1.5em]
  \item \textbf{Homogen}, spalten-major (wie Fortran), Index \texttt{[i,j]}
  \item Lineare Algebra: \texttt{t(M)} (Transpose), \texttt{M \%*\% N} (Matrixmultiplikation)
  \item Vorteil: Numerisch stabil, schnell; Nachteil: nur ein Datentyp
\end{itemize}

\textbf{Arrays -- Allgemeine multidimensionale Strukturen:}\vspace{2pt}
\begin{itemize}[leftmargin=1.5em]
  \item $n$-dimensionale, homogene Struktur (Verallgemeinerung von Matrix)
  \item Index \texttt{arr[i1, i2, ..., in]}; für Teilmengen: \texttt{arr[1:5, , 3]}
  \item Einsatz: Zeitreihenanalyse mit mehreren Variablen, 3D-Daten (z.B. Neuronale Netze Pre-Processing)
\end{itemize}

\definition{\textbf{Data Frame}: Tabellenstruktur mit Spalten verschiedener Datentypen (heterogen wie SQL-Tabelle). Spalten sind Vektoren gleicher Länge, indexierbar via \texttt{\$}-Operator oder \texttt{[,]}.}

\textbf{Data Frame Eigenschaften:}\vspace{2pt}
\begin{itemize}[leftmargin=1.5em]
  \item \textbf{Heterogen}: Spalten können \texttt{numeric, character, factor, logical} sein
  \item \textbf{Effizient für Statistik}: \texttt{lm(y \textasciitilde\ x1 + x2, data=df)}
  \item Vorteil: Flexibel, SQL-ähnlich; Nachteil: langsamer als Matrix bei Numerik
  \item Moderne Alternative: \texttt{tibble} (verbesserter DF), \texttt{data.table} (schnell)
\end{itemize}

\textbf{Listen -- Generische Vektoren:}\vspace{2pt}
\begin{itemize}[leftmargin=1.5em]
  \item Heterogen: Elemente können beliebige Typen haben (auch andere Listen)
  \item Index via \texttt{list[[i]]}, \texttt{list\$name}
  \item Einsatz: Funktionsrückgaben (mehrere Objekte), komplexe Datenstrukturen
  \item Nachteil: Langsamer als homogene Strukturen
\end{itemize}

\begin{table}[h!]
\centering
\rowcolors{2}{tablerowB}{tablerowA}
\begin{tabularx}{\linewidth}{L{2.2cm} C{1.8cm} L{4.8cm} L{3.2cm}}
\rowcolor{tablehead}
\color{white}\textbf{Struktur} &
\color{white}\textbf{Dim} &
\color{white}\textbf{Eigenschaft} &
\color{white}\textbf{Typischer Use-Case} \\
\midrule
Vektor & 1D & Homogen, vektorisiert & Numerische Sequenzen \\
Matrix & 2D & Homogen, spalten-major & Lineare Algebra \\
Array & nD & Homogen, beliebige Dim & Multidimensionale Daten \\
Data Frame & 2D & Heterogen (Spalten) & Statistische Analyse \\
Liste & nD & Heterogen (rekursiv) & Komplexe Strukturen \\
\end{tabularx}
\end{table}

\newpage

\subsection{Objektorientierte Programmierung in R}

\subsubsection{Überblick und Besonderheit}

R ist primär eine funktionale Sprache -- OOP ist ein nachträglich hinzugefügtes Konzept. Deshalb gibt es \textbf{mehrere, teils konkurrierende OOP-Systeme}, die historisch gewachsen sind. Die vier wichtigsten sind: \textbf{S3}, \textbf{S4}, \textbf{R5 (Reference Classes)} und \textbf{R6}.

\subsubsection{S3 -- Informelles Dispatch-System}

S3 ist das älteste und am häufigsten verwendete OOP-System in R. Es basiert auf \textbf{generischen Funktionen und Method Dispatch}: Eine generische Funktion (z.\,B.\ \texttt{print}, \texttt{summary}, \texttt{plot}) ruft die passende Methode auf Basis des \texttt{class}-Attributs des ersten Arguments auf.

\textbf{Method Dispatch:} Wenn \texttt{speak(fido)} aufgerufen wird, sucht R nach \texttt{speak.Dog}, dann \texttt{speak.default}. Die Vererbung wird über den \texttt{class}-Vektor mit mehreren Einträgen realisiert.

\textbf{Vorteile:} Extrem leichtgewichtig; keine formale Definition nötig; kompatibel mit den meisten Base-R-Paketen. \\
\textbf{Nachteile:} Keine formale Klassenvalidierung; fehleranfällig bei großen Projekten.

\subsubsection{S4 -- Formales Klassen-System}

S4 ist ein formales OOP-System mit expliziter Klassendefinition, Slot-Validierung und Mehrfachvererbung. Es wird vor allem in Bioinformatik-Paketen (Bioconductor) eingesetzt.

\textbf{Schlüsselunterschiede zu S3:}
\begin{itemize}
  \item \textbf{Formale Deklaration} via \texttt{setClass} -- Slots (Felder) und ihre Typen sind explizit definiert
  \item \textbf{Slot-Zugriff} via \texttt{@} (nicht \texttt{\$})
  \item \textbf{Validierung} via \texttt{validity}-Argument möglich
  \item \textbf{Mehrfachvererbung} via \texttt{contains}
  \item \textbf{Multiple Dispatch}: \texttt{setMethod} kann auf der Kombination mehrerer Argument-Typen dispatchen
\end{itemize}

\textbf{Vor- und Nachteile:} Formal, sicher und gut dokumentierbar, aber verbos und deutlich komplexer als S3.

\subsubsection{R5 / Reference Classes -- Mutable OOP}

Reference Classes (auch R5) bringen \textbf{echte Referenzsemantik} nach R -- ähnlich wie Klassen in Python oder Java. Objekte sind mutable: Methoden können Felder in-place ändern (kein Copy-on-Modify).

\textbf{Vor- und Nachteile:} Referenzsemantik kann in manchen Situationen (z.\,B.\ Zustandsautomaten, Caches) eleganter sein. Allerdings ungewohnt für R-Programmierer; Debugging schwieriger; kaum in Standard-Paketen verwendet.

\subsubsection{R6 -- Modernes, schlankes Referenz-System}

Das Paket \texttt{R6} bietet eine moderne Alternative zu R5 mit saubererer Syntax, echter privater/öffentlicher Felddefinition und besserer Performance.

\subsubsection{Vergleichstabelle der OOP-Systeme}

\begin{table}[h!]
\centering
\rowcolors{2}{tablerowB}{tablerowA}
\begin{tabularx}{\linewidth}{L{2.5cm} C{2cm} C{2cm} C{2cm} C{2cm}}
\rowcolor{tablehead}
\color{white}\textbf{Merkmal} &
\color{white}\textbf{S3} &
\color{white}\textbf{S4} &
\color{white}\textbf{R5} &
\color{white}\textbf{R6} \\
\midrule
Formale Deklaration & -- & \checkmark & \checkmark & \checkmark \\
Referenzsemantik    & -- & --         & \checkmark & \checkmark \\
Kapselung (private) & -- & --         & eingeschränkt & \checkmark \\
Multiple Dispatch   & -- & \checkmark & --         & -- \\
Komplexität         & gering & hoch  & mittel     & mittel \\
Performance         & gut & gut       & ok         & gut \\
Verbreitung         & sehr hoch & hoch & gering  & mittel \\
\end{tabularx}
\end{table}

\subsubsection{Empfehlungen}

\begin{itemize}
  \item \textbf{S3} für einfache Pakete und schnelle Prototypen -- die Default-Wahl in R
  \item \textbf{S4} wenn formale Typsicherheit und Multiple Dispatch benötigt werden (z.\,B.\ Bioconductor)
  \item \textbf{R6} wenn Zustandsverwaltung mit echter Kapselung gewünscht ist -- am nächsten an Python/Java-OOP
  \item \textbf{R5} heute weitgehend durch R6 ersetzt
\end{itemize}

\newpage

\subsection{Tidyverse -- Philosophie \& Funktionsweise}

\subsubsection{Was ist das Tidyverse?}

Das Tidyverse ist ein \textbf{kohärentes Ökosystem von R-Paketen}, das von Hadley Wickham und dem RStudio (heute Posit) Team entwickelt wurde. Es teilt eine gemeinsame Design-Philosophie, konsistente API und Datenstruktur. Kernpakete sind: \texttt{ggplot2}, \texttt{dplyr}, \texttt{tidyr}, \texttt{readr}, \texttt{purrr}, \texttt{tibble}, \texttt{stringr}, \texttt{forcats}.

\subsubsection{Die Kernphilosophie: Tidy Data}

Das Tidyverse basiert auf dem Konzept \textbf{Tidy Data} (Hadley Wickham, 2014):
\begin{enumerate}
  \item Jede \textbf{Variable} bildet eine \textbf{Spalte}
  \item Jede \textbf{Beobachtung} bildet eine \textbf{Zeile}
  \item Jeder \textbf{Wert} steht in genau \textbf{einer Zelle}
\end{enumerate}

Tidy Data ist das standardisierte Eingabeformat für alle Tidyverse-Funktionen und eliminiert das häufige Reshaping, das in Base-R notwendig wäre.

\subsubsection{Das Tibble -- moderner Data Frame}

\texttt{tibble} ist eine überarbeitete Version des Data Frames mit besserem Druck-Verhalten, strikterer Subsetting-Semantik (kein automatisches Type-Dropping) und schnellerer Erstellung.

\subsubsection{dplyr -- Data Manipulation}

\texttt{dplyr} stellt fünf \textbf{Kern-Verben} bereit, die die häufigsten Datentransformationen abdecken:
\begin{itemize}
  \item \textbf{filter():} Zeilen filtern
  \item \textbf{select():} Spalten auswählen
  \item \textbf{mutate():} Neue Spalte berechnen
  \item \textbf{arrange():} Sortieren
  \item \textbf{summarise():} Aggregieren
\end{itemize}

\textbf{Der Pipe-Operator:} Zentrales Stilelement ist der Pipe-Operator \texttt{|>} (seit R 4.1) oder \texttt{\%>\%} (aus \texttt{magrittr}). Der Pipe übergibt das Ergebnis der linken Seite als erstes Argument der rechten Seite -- eine links-nach-rechts-lesbare Kette von Transformationen.

\textbf{Gruppierung und Aggregation:} \texttt{group\_by()} + \texttt{summarise()} ermöglichen komplexe Aggregationen pro Gruppe.

\subsubsection{tidyr -- Daten in Tidy Format}

\texttt{tidyr} bietet Funktionen zum Reshaping:
\begin{itemize}
  \item \texttt{pivot\_longer():} Wide $\to$ Long Format (häufig nötig für ggplot2)
  \item \texttt{pivot\_wider():} Long $\to$ Wide Format
  \item \texttt{unnest():} Verschachtelte Daten expandieren
\end{itemize}

\subsubsection{ggplot2 -- Grafiken nach dem Grammar of Graphics}

\texttt{ggplot2} implementiert die \textbf{Grammar of Graphics}: Ein Plot wird aus Schichten (Layers) aufgebaut, bestehend aus Daten (\texttt{aes}thetics) und geometrischen Objekten (\texttt{geom}s). Das ermöglicht eine deklarative, konsistente Beschreibung beliebig komplexer Grafiken.

In Base-R würde dieselbe Grafik \texttt{plot()}, \texttt{points()}, \texttt{lines()}, \texttt{legend()} etc.\ erfordern -- jedes Mal mit manueller Koordination.

\subsubsection{purrr -- Funktionale Programmierung}

\texttt{purrr} bietet eine konsistente Familie von \texttt{map\_*}-Funktionen als typsichere Alternative zu \texttt{lapply}/\texttt{sapply}:
\begin{itemize}
  \item \texttt{map()} gibt immer eine Liste zur\u00fcck
  \item \texttt{map\_dbl():} Ergebnis ist garantiert ein double-Vektor
  \item \texttt{map2():} Zwei Listen parallel verarbeiten
\end{itemize}

Das \texttt{~}-Syntax erzeugt eine anonyme Funktion (\texttt{.x} referenziert das aktuelle Element).

\subsubsection{Vorteile gegenüber plain R}

\begin{table}[h!]
\centering
\rowcolors{2}{tablerowB}{tablerowA}
\begin{tabularx}{\linewidth}{L{3.5cm} L{5.2cm} L{5.2cm}}
\rowcolor{tablehead}
\color{white}\textbf{Aspekt} &
\color{white}\textbf{Base R} &
\color{white}\textbf{Tidyverse} \\
\midrule
Lesbarkeit       & Verschachtelter Code & Lineare Pipe-Kette \\
Konsistenz       & Inkonsistente Funktionsnamen & Einheitliche Verb-API \\
Datenstruktur    & \texttt{data.frame} mit Fallen & \texttt{tibble} (strikt, sicher) \\
Grafik           & \texttt{plot()} (imperativ) & \texttt{ggplot2} (deklarativ) \\
String-Ops       & \texttt{grep}, \texttt{gsub} (inkonsistent) & \texttt{stringr} (konsistent) \\
Iteration        & \texttt{lapply}/\texttt{sapply} & \texttt{purrr::map\_*} (typsicher) \\
Faktor-Handling  & Historisch fehleranfällig & \texttt{forcats} \\
\end{tabularx}
\end{table}

\subsubsection{Nachteile und Kritik}

\begin{itemize}
  \item \textbf{Abhängigkeiten:} Das Tidyverse bringt viele Pakete mit -- Dependency-Management kann aufwändig sein
  \item \textbf{Performance:} \texttt{dplyr} kann bei sehr großen Datensätzen langsamer sein als spezialisierte Pakete wie \texttt{data.table}
  \item \textbf{Non-standard Evaluation (NSE):} \texttt{dplyr} nutzt NSE -- elegant, aber erschwert programmatische Nutzung in Funktionen
  \item \textbf{Versionsstabilität:} Das Tidyverse entwickelt sich schnell; Breaking Changes erfordern Code-Aktualisierung
\end{itemize}

\newpage

\subsection{Wichtige Dateiformate}

\subsubsection{Überblick}

Dateiformate lassen sich nach mehreren Kriterien klassifizieren:
\begin{itemize}
  \item \textbf{Text vs. Binär:} Text (CSV, JSON, YAML, XML) ist human-readable; Binär (HDF5, Parquet, Pickle) ist kompakt
  \item \textbf{Hierarchisch vs. tabellarisch:} XML, JSON, YAML sind hierarchisch; CSV, Parquet sind tabellarisch
  \item \textbf{Schema vs. schemalos:} XML unterstützt formale Schemata (XSD); JSON Schema ist optional; CSV und YAML sind vollständig schemalos
\end{itemize}

\subsubsection{CSV -- Comma-Separated Values}

CSV ist das universellste Format für tabellarische Daten. Jede Zeile = ein Datensatz; Felder werden durch Trennzeichen (meist Komma, in DACH oft Semikolon) getrennt.

\textbf{Vorteile:} Universell kompatibel; menschenlesbar; kein Parser-Overhead. \\
\textbf{Nachteile:} Kein Datentyp-System; keine Unterstützung für Hierarchien; Encoding-Probleme; sehr große CSV-Dateien sind langsam.

\subsubsection{XML -- eXtensible Markup Language}

XML ist hierarchisch und textbasiert. Daten werden in \textbf{Tags} eingeschlossen; Tags können beliebig verschachtelt werden.

\textbf{Vorteile:} Formale Validierung via XSD; weit verbreitet im Unternehmensumfeld; selbstbeschreibend. \\
\textbf{Nachteile:} Sehr verbose; schwerer lesbar als JSON; teurer Parsing-Overhead.

\subsubsection{JSON -- JavaScript Object Notation}

JSON ist das de-facto-Standardformat für Web-APIs. Es kennt sechs Datentypen: \texttt{object}, \texttt{array}, \texttt{string}, \texttt{number}, \texttt{boolean}, \texttt{null}.

\textbf{Warum JSON XML verdrängt hat:} JSON ist direkter auf Programmiersprachenstrukturen (Dicts, Arrays) abbildbar; weniger verbose; leichter in JavaScript zu verarbeiten. REST-APIs verwenden fast ausschließlich JSON.

\textbf{Einschränkungen:} Keine Kommentare, keine nativen Datumstypen, keine Unterscheidung zwischen Integer und Float auf Typsystemebene.

\subsubsection{YAML -- YAML Ain't Markup Language}

YAML ist eine Obermenge von JSON und wurde für maximale menschliche Lesbarkeit optimiert. Einrückung definiert die Hierarchie; Anführungszeichen sind optional.

\textbf{Typische Anwendungsfälle:} Konfigurationsdateien (Docker Compose, Kubernetes, GitHub Actions, Ansible, CI/CD-Pipelines).

\textbf{Gefahren:} YAML-Parsing ist komplex und fehleranfällig. \\
\textbf{Sicherheit:} Immer \texttt{yaml.safe\_load()} verwenden, nie \texttt{yaml.load()} aus unbekannten Quellen!

\subsubsection{HDF5 -- Hierarchical Data Format}

HDF5 ist ein binäres, hierarchisches Dateiformat. Eine HDF5-Datei ist wie ein Dateisystem: es gibt \textbf{Groups} (Verzeichnisse) und \textbf{Datasets} (n-dimensionale Arrays).

\textbf{Warum HDF5 für ML?} Lazy Loading: Nur angeforderte Daten werden vom Disk geladen. Keras speichert Modellgewichte in HDF5. Numpy-Arrays werden direkt gespeichert. Kompression reduziert Größe.

\textbf{Nachteile:} Nicht human-readable; Parallelschreiben eingeschränkt; weniger flexibel als Parquet.

\subsubsection{Parquet}

Apache Parquet ist spaltenorientiert und aus dem Hadoop-Ökosystem. Alle Werte einer Spalte werden zusammen gespeichert -- extrem effiziente Lesezugriffe bei wenigen Spalten (OLAP-Workload). Integrierte Kompression und Typsicherheit.

\textbf{Standard in:} Apache Spark, AWS Athena, Google BigQuery, DuckDB.

\subsubsection{Weitere Formate: Pickle, NetCDF}

\textbf{Pickle:} Python-spezifisches Binärformat. Praktisch, aber \textbf{sicherheitskritisch} -- nie aus unbekannten Quellen laden!

\textbf{NetCDF:} Standard in Klimawissenschaft und Meteorologie für multidimensionale Array-Daten.

\newpage

\subsection{git}

\subsubsection{Was ist Git?}

Git ist ein \textbf{verteiltes Versionskontrollsystem} (DVCS), ursprünglich 2005 von Linus Torvalds geschrieben. Im Gegensatz zu zentralisierten Systemen (SVN, CVS) hat jeder Entwickler eine \textbf{vollständige Kopie} des Repositories -- inkl. kompletter Historie -- auf seinem Rechner.

\subsubsection{Kernkonzepte}

\textbf{Das Git-Objektmodell:} Git speichert \textbf{Snapshots}, nicht Deltas. Vier Objekttypen im \texttt{.git/objects/}:
\begin{itemize}
  \item \textbf{Blob:} Dateiinhalt (beliebige Binärdaten)
  \item \textbf{Tree:} Verzeichnisstruktur (Liste von Blobs und Sub-Trees)
  \item \textbf{Commit:} Snapshot eines Projektzustands mit Author, Zeitstempel, Parent-Commit
  \item \textbf{Tag:} Benannte Referenz auf einen Commit (z.B. \texttt{v1.0.0})
\end{itemize}

Alle Objekte werden durch SHA-1-Hash (40-stellige Hex-Zeichenkette) identifiziert. Jede Änderung ergibt einen neuen Hash -- Historia ist kryptographisch manipulationssicher.

\textbf{Die drei Bereiche:}
\begin{itemize}
  \item \textbf{Working Directory:} Deine Dateien auf der Festplatte
  \item \textbf{Staging Area:} Zwischenspeicher für \texttt{git add} -- man wählt aus, was in den nächsten Commit geht
  \item \textbf{Repository (.git/):} Die lokale Datenbank mit aller Historie
\end{itemize}

\subsubsection{Branches und Merging}

Ein \textbf{Branch} ist nur ein Zeiger auf einen Commit. Erstellen ist eine $\mathcal{O}(1)$-Operation -- kein Kopieren von Dateien.

\textbf{Merge:} Zwei Commit-Linien zusammenführen. Zwei Typen:
\begin{itemize}
  \item \textbf{Fast-Forward:} Wenn eine Branch einfach weiter vorne ist, wird der Zeiger verschoben
  \item \textbf{Rebase:} Alternative -- schreibt Commits so um, als wären sie auf aktuellem Stand entstanden. Erzeugt lineare Geschichte, verändert aber SHA-1-Hashes (daher nicht auf shared Branches!)
\end{itemize}

\subsubsection{Basisbefehle}

\textbf{Repository einrichten:}
\begin{itemize}
  \item \texttt{git init} -- Neues lokales Repo erstellen
  \item \texttt{git clone <url>} -- Bestehendes Repo klonen
\end{itemize}

\textbf{Status und Unterschiede:}
\begin{itemize}
  \item \texttt{git status} -- Zustand anzeigen
  \item \texttt{git diff} -- Änderungen im Working Dir
  \item \texttt{git log --oneline --graph} -- Geschichte mit Branch-Grafik
\end{itemize}

\textbf{Änderungen erfassen:}
\begin{itemize}
  \item \texttt{git add <datei>} -- Datei stagen
  \item \texttt{git add -p} -- Interaktiv: nur bestimmte Hunks
  \item \texttt{git commit -m "Nachricht"} -- Commit mit Message
  \item \texttt{git commit --amend} -- Letzten Commit korrigieren
\end{itemize}

\textbf{Branches:}
\begin{itemize}
  \item \texttt{git branch} -- Branches auflisten
  \item \texttt{git switch -c <branch>} -- Branch erstellen und wechseln
  \item \texttt{git merge <branch>} -- Mergen
  \item \texttt{git rebase main} -- Rebasing
\end{itemize}

\textbf{Remote-Operationen:}
\begin{itemize}
  \item \texttt{git fetch origin} -- Änderungen vom Remote holen (OHNE merge)
  \item \texttt{git pull origin main} -- fetch + merge
  \item \texttt{git push origin main} -- Commits zum Remote senden
\end{itemize}

\textbf{Rückgängig machen:}
\begin{itemize}
  \item \texttt{git restore <datei>} -- Working-Dir-Änderung verwerfen
  \item \texttt{git revert <hash>} -- Neuen Commit erstellen, der Commit rückgängig macht
  \item \texttt{git reset --hard HEAD\textasciitilde 1} -- Letzten Commit rückgängig (DATEN VERLOREN!)
  \item \texttt{git reflog} -- Rettungsanker nach \texttt{reset --hard}
\end{itemize}

\subsubsection{Workflow und Best Practices}

\textbf{Typischer Feature Branch Workflow:}
\begin{enumerate}
  \item \texttt{git switch main \&\& git pull origin main}
  \item \texttt{git switch -c feature/my-feature}
  \item Entwickeln und committen
  \item \texttt{git push -u origin feature/my-feature}
  \item Pull Request erstellen für Code Review
  \item Nach Merge: \texttt{git branch -d feature/my-feature}
\end{enumerate}

\textbf{.gitignore:} Definiert, welche Dateien Git ignorieren soll. Typische Einträge:
\begin{itemize}
  \item \texttt{\_\_pycache\_\_/}, \texttt{*.pyc}, \texttt{.venv/} (Python)
  \item \texttt{data/raw/}, \texttt{*.csv}, \texttt{*.h5}, \texttt{*.pkl} (Daten -- nutze Git LFS oder DVC!)
  \item \texttt{.vscode/}, \texttt{.idea/} (IDE)
\end{itemize}

\newpage